\subsubsection{Método}\label{sec:transicion-metodo}La fase de Transición de DAD se adapta como una semi-transición, orientada a evaluar y entregar académicamente una primera versión funcional, sin contemplar su despliegue e integración definitiva en una institución encargada de atender emergencias. La aplicación se implementará para el sector de Aguas Claras, en los Cerros Orientales de Bogotá, a partir de la información geoespacial, meteorológica y de vegetación disponible durante el proyecto. Esta primera aproximación podría ampliarse a otros sectores o territorios si se cuenta con datos adecuados, especialmente de vegetación y combustibles.

La finalidad principal será evaluar con comandantes o coordinadores de emergencias si la aplicación responde a la necesidad identificada de consultar e interpretar información sobre la posible propagación de un incendio forestal. La participación en la evaluación no constituye aceptación institucional ni autorización para el uso real.
